\documentclass[../NDCNNAI_main.tex]{subfiles}
\graphicspath{{\subfix{../fig/}}}

\begin{document}
\subsection{Универсальная аппроксимация операторов}
\begin{theorem}[Чен и Чен, теорема 5]
    \label{thm:chen_chen_oper}
    Пусть \( X \) --- банахово пространство, \( K_1 \subset X \) компактно, \( K_2 \subset \mathbb{R}^n \) компактно, \( V \subset C(K_1) \) компактно и \( G \colon V \to C(K_2) \) непрерывно. Пусть \( g \) является функцией Таубер–Винера. Тогда для любого \( \varepsilon > 0 \) существуют целые числа \( m, M, N \), точки выборки \( x_1, \ldots, x_m \in K_1 \), выходные частоты \( w_1, \ldots, w_N \in \mathbb{R}^n \) и вещественные параметры \( \{\theta_{ij}, \gamma_i\}_{i \leq M, j \leq m} \) и \( \{\alpha_{ik}, \beta_k\}_{i \leq M, k \leq N} \) такие, что
\[
\sup_{u \in V} \sup_{y \in K_2} \left| G(u)(y) - \sum_{k=1}^{N} \left( \sum_{i=1}^{M} \alpha_{ik} g \left( \sum_{j=1}^{m} \theta_{ij} u(x_j) + \gamma_i \right) \right) g(w_k \cdot y + \beta_k) \right| \leq \varepsilon.
\]
\end{theorem}

\textbf{Эквивалентная формулировка:} двухслойное сепарабельное (<<внешнее произведение>>) разложение по \( g \)-атомам над входными выборками и выходным значением дает равномерную аппроксимацию оператора.


\begin{proof}
\textbf{Шаг 0 (компактный образ).} Поскольку \( V \) компактно в \( C(K_1) \) и \( G \) непрерывно, образ  
\[
U \coloneq  G(V) \subset C(K_2)
\]
является компактным.

\textbf{Шаг 1 (применение теоремы об равномерной УА на \( U \subset C(K_2) \)).} По теореме с \( K = K_2 \) и компактным семейством \( U \), для \( \varepsilon_1 > 0 \) существуют \( N \in \mathbb{N} \), параметры \( (w_k, \beta_k)_{k=1}^{N} \) (независящие от \( u \)) и непрерывные линейные функционалы \( L_k \colon U \to \mathbb{R} \) такие, что для всех \( u \in V \),
\[
\left\| G(u)(\cdot) - \sum_{k=1}^{N} L_k(G(u)) g(w_k\cdot(\cdot) + \beta_k) \right\|_{C(K_2)} \leq \varepsilon_1.
\]
Таким образом, <<коэффициенты как функции \( u \)>> в точности являются скалярами \( L_k(G(u)) \).

\textbf{Шаг 2: сведение \( L_k \circ G \colon V \rightarrow \mathbb{R} \) к \( g \)-сети на основе выборок}
Для каждого фиксированного \( k \) отображение \( F_k \coloneq  L_k \circ G \colon V \rightarrow \mathbb{R} \) непрерывно (композиция непрерывных отображений на компактных множествах). Применим теорему об УА для функционалов к каждому \( F_k \) с точностью \( \varepsilon_{2,k} > 0 \). Получим:

\begin{itemize}
\item конечное множество выборок \( \{x_j\}_{j=1}^m \subset K_1 \) и \textit{внутренние} параметры \( \{\theta_{ij}, \gamma_i\}_{i=1}^M \), \textit{единые} для (конечного) набора \( \{F_k\}_{k=1}^N \) (объединим точки для каждого \(F_k \)),
\item внешние коэффициенты \( \{\alpha_{ik}\}_{i=1,\ldots,M} \) (зависящие от \( k \) и \( \{F_k\} \), но не от \( u \)),
такие, что для всех \( u \in V \) и всех \( k = 1, \ldots, N \),

\[
\left| F_k(u) - \sum_{i=1}^M \alpha_{ik} g \left( \sum_{j=1}^m \theta_{ij} u(x_j) + \gamma_i \right) \right| \leq \varepsilon_{2,k}.
\]
\end{itemize}
(Использовали, что теорема об УА для функционалов определяет внутренние параметры независимо от функционала и меняет только внешние коэффициенты.)

\textbf{Шаг 3: композиция и отслеживание \( \varepsilon \)}
Подставим аппроксимацию из шага 2 в разложение шага 1. Для любых \( u \in V \) и \( y \in K_2 \):

\begin{multline*}
\left|
G(u)(y) - \sum_{k=1}^{N} \left( \sum_{i=1}^{M} \alpha_{ik} g \left( \sum_{j=1}^{m} \theta_{ij} u(x_j) + \gamma_i \right) \right) g(w_k \cdot y + \beta_k)
\right| \leq\\
\leq 
\underbrace{\left\| G(u) - \sum_{k=1}^{N} L_k (G(u)) g(w_k \cdot + \beta_k) \right\|_{C(K_2)}}_{\leq \varepsilon_1} + \sum_{k=1}^{N} \underbrace{\big | F_k (u) - \sum_{i=1}^{M} \alpha_{ik} g(\cdots) \big|}_{\leq \varepsilon_{2,k}} \left\| g(w_k \cdot + \beta_k) \right\|_{C(K_2)}.
\end{multline*}

Положим \( B_k \coloneq  \left\| g(w_k \cdot + \beta_k) \right\|_{C(K_2)} < \infty \) (непрерывность ridge-функции \( g \) на компакте \( K_2 \)). Выберем

\[
\varepsilon_1 \coloneq  \varepsilon/3, \quad \varepsilon_{2,k} \coloneq  \frac{\varepsilon}{3 \sum_{\ell=1}^{N} B_\ell}.
\]
Тогда общая ошибка \( \leq \varepsilon/3 + \sum_{k} \varepsilon_{2,k} B_k \leq \varepsilon/3 + \varepsilon/3 < \varepsilon \), равномерно по \( u \in V \) и \( y \in K_2 \). Это даёт заявленную сепарабельную архитектуру.
\end{proof}

Отметим, почему получилось доказать теорему, нам удалось \textbf{сохранить структуру}, потому что:
\begin{itemize}
    \item Выходные атомы \( (w_k, \beta_k) \) получены из теоремы о равномерной УА на компактном образе \( U = G(V) \) и не зависят от \( u \).

    \item Внутренние точки выборки \( x_j \) и внутренние параметры \( (\theta_{ij}, \gamma_i) \) выбираются один раз (для всех \( k \)) применением теоремы об УА для функционалов к конечному набору \( \{F_k\}_{k=1}^N \) с последующим объединением.

    \item Только скаляры \( \alpha_{ik} \) (коэффициенты смещения) зависят от того, какой коэффициентный функционал \( F_k \) аппроксимируется; они являются константами в конечном операторе и не зависят от \( u \) или \( y \).
\end{itemize}
\textbf{Равномерность} обеспечивает следующее наблюдение.
Каждый шаг равномерен на компактных множествах: \( U = G(V) \) компактно, Теорема B равномерна на \( U \), а Теорема C равномерна на \( V \). Следовательно, итоговая оценка является \textit{равномерной} как по \( u \in V \), так и по \( y \in K_2 \).

\begin{remarks}
    \leavevmode
    \begin{itemize}
    \item  \textbf{Динамическая взгляд.} Когда на \( K_1/K_2 \) определено действие сдвига/потока, ridge-атомы \( g(w \cdot y + \beta) \) описывают переносы и растяжения в выходном пространстве, тогда как внутренняя операция смешения выборок реализует конечную «память» входа \( u \) --- в точности ту структуручто применяется при идентификации инвариантных по времени/пространству систем.

    \item  \textbf{Векторнозначные выходы.} Для \( G \colon V \to C(K_2; \mathbb{R}^d) \) применяем теорему к каждой компоненте и объединяем полученные разложения.

    \item  \textbf{Альтернативный набросок доказательства.} Можно также зафиксировать \( y \in K_2 \), применить приблизить функционал вычисления \( u \mapsto G(u)(y) \) равномерно по \( y \) (используя теорему о равномерной УА на \( K_2 \) для разделения параметров по \( y \)), что даёт ту же сепарабельную форму.
    \end{itemize}    
\end{remarks}

%%%%%%%%%%%%%%%%%%%%%%%%%%%%%%%%%%%%%%%%%%%%%%%


\subsection{Взгляд с точки зрения динамических систем и сдвиги}

\textbf{Постановка.} Пусть (топологическая) группа \( G \) действует непрерывно на компактных множествах \( K_1, K_2 \): \( G \times K_\ell \to K_\ell \), \((g, x) \mapsto g \cdot x\). Индуцированное действие на \( C(K_\ell) \) имеет вид

\[
(g \cdot u)(x) \coloneq  u(g^{-1} \cdot x), \quad u \in C(K_\ell).
\]

Отображение \( f \colon V \subset C(K_1) \to \mathbb{R} \) является \( \rho \)-эквивариантным (\( \rho \colon G \to \mathbb{C}\)), если \( f(g \cdot u) = \rho(g)f(u) \).  
Отображение \( G \colon V \to C(K_2) \) является эквивариантным, если \( G(g \cdot u) = g \cdot G(u) \).

\textbf{TW-устойчивость атомов.} Ridge-атомы \( x \mapsto g(w \cdot x + \theta) \) устойчивы относительно базовых действий:

\[
\underbrace{T_{a} g(w \cdot x + \theta)}_{\text{сдвиг по } x} = g(w \cdot x + (\theta + w \cdot a)),
\quad
\underbrace{D_{\lambda} g(w \cdot x + \theta)}_{\text{растяжение}} = g((\lambda w) \cdot x + \theta).
\]
Таким образом, семейство атомов замкнуто относительно действий сдвига/растяжения.


\begin{theorem}[ДС-A, эквивариантная УА для функционалов]

Пусть \( G \) действует на метрическом компакте \( K \) и \( V \subset C(K) \) является непустым компактным \( G \)-инвариантным множеством.  
Пусть \( f \colon V \to \mathbb{C} \) непрерывно и \( \rho \)-эквивариантно: \( f(g \cdot u) = \rho(g)f(u) \) для всех \( g \in G \).  
Предположим, что \( g \) (активация) является ТВ-функцией. Тогда для любого \( \varepsilon > 0 \) существуют

\begin{itemize}
    \item точки выборки \( x_1, \ldots, x_m \in K \) и внутренние параметры \( \{\theta_{ij}, \beta_i\} \) (не зависящие от \( u \)),
\item внешние коэффициенты \( c_i \in \mathbb{C} \),
\end{itemize}
такие, что <<выборочная нейросеть>>

\[
F(u) = \sum_{i=1}^{N} c_i \, g \left( \sum_{j=1}^{m} \theta_{ij} u(x_j) + \beta_i \right)
\]
удовлетворяет $\sup_{u \in V} |F(u) - f(u)| < \varepsilon$, и её можно выбрать \( \rho \)-эквивариантной.
\end{theorem}

\begin{proof} (i) Применить теорему об УА для функционалов к \( f \) для получения равномерной аппроксимации на основе выборок на \( V \).
(ii) Обеспечить \( \rho \)-эквивариантность путём усреднения коэффициентов по группе:  
\[
\tilde{F}(u) \coloneq  \int_G \rho(g)^{-1} F(g \cdot u) \, d\mu(g)
\]  
(мера Хаара).  
Тогда \(\tilde{F}\) является \( \rho \)-эквивариантной, зависит от тех же внутренних параметров, и ошибка не возрастает.
\end{proof}


\begin{theorem}[ДС-B, биэквивариантная УА для операторов]

Пусть \( G \) действует на компактах \( K_1, K_2 \), \( V \subset C(K_1) \) компактно и \( G \)-инвариантно, и \( G : V \to C(K_2) \) непрерывен и эквивариантен:
\(
G(h \cdot u) = h \cdot G(u) \text{ для всех } h \in G.
\)

Предположим, что активация \( g \) является TW-функцией. Тогда для любого \( \varepsilon > 0 \) существуют
\begin{itemize}
\item входные точки выборки \( x_1, \ldots, x_m \in K_1 \) и внутренние параметры \( (\theta_{ij}, \gamma_i) \),
\item выходные атомы \( y \mapsto g(w_k \cdot y + \beta_k), \, k = 1, \ldots, N \) (общие для всех \( u \)),
\item коэффициенты смешения \( \alpha_{ik} \),
\end{itemize}
такие, что сепарабельная архитектура
\[
\mathcal{G}(u)(y) \coloneq  \sum_{k=1}^{N} \left( \sum_{i=1}^{M} \alpha_{ik} g \left( \sum_{j=1}^{m} \theta_{ij} u(x_j) + \gamma_i \right) \right) g(w_k \cdot y + \beta_k)
\]
является эквивариантной и \( \sup_{u \in V} \|\mathcal{G}(u) - G(u)\|_{C(K_2)} < \varepsilon \).
\end{theorem}

\begin{proof}
Применить теорему об УА для операторов, чтобы получить сепарабельную аппроксимацию, затем сделать её эквивариантной путём усреднения по \( G \) на выходном слое:  
\[
\tilde{\mathcal{G}}(u) \coloneq  \int_G h \cdot \mathcal{G}(h^{-1} \cdot u) d\mu(h).
\]  
TW-устойчивость атомов гарантирует, что усреднение перепараметризует атомы внутри того же ridge-семейства.   
\end{proof}


\subsection{Устойчивость относительно потоков, робастность аппроксимаций}

Пусть \( \{\Phi_t\}_{t \in \mathbb{R}} \) --- непрерывный поток на \( K \) (например, сдвиги), с индуцированным действием на \( C(K) \). Предположим, что \( V \subset C(K) \) компактно и инвариантно относительно \( \Phi_t \) и \( f \colon V \to \mathbb{R} \) (или \( G \colon V \to C(K) \)) непрерывен и эквивариантен/инвариантен относительно \( \Phi_t \).

Отметим \textbf{свойство.} Усреднённые сети из ДС-A/ДС-B удовлетворяют  
\[
F(\Phi_t \cdot u) = \rho(t) F(u), \quad \mathcal{G}(\Phi_t \cdot u) = \Phi_t \cdot \mathcal{G}(u), \quad \forall t \in \mathbb{R},
\]
и оценки ошибки равномерны по \( t \) на компактных временных интервалах.

\textbf{Причина} этого явления:
\begin{enumerate}[label=(\roman*)]
    \item Усреднение по группе/потоку коммутирует с sup-нормой и сохраняет равномерные оценки.
    \item Ridge-атомы \( g(w \cdot x + \theta) \) остаются в том же параметрическом семействе при действии \( \Phi_t \) (сдвиг/растяжение перепараметризуют \( (w, \theta) \)), поэтому скрытый слой является устойчивым к действию.
\end{enumerate}

Как итог можно привести динамические формулировки теорем этой лекции.
\begin{itemize}
\item \textbf{Теорема A как динамический критерий}. ТВ \( \Leftrightarrow \) «неполиномиальность» означает: ridge-атомы, порождённые действиями сдвига/растяжения, образуют динамически устойчивый словарь.

\item \textbf{Теорема B как равномерная УА с устойчивостью относительно действия}. На компактном \( U \subset C(K) \) аппроксимируем с помощью фиксированного скрытого слоя, чьи атомы устойчивы относительно действия на \( K \); варьируются только линейные внешние функционалы.

\item \textbf{Теоремы C/D как эквивариантные УА}. После усреднения по группе/потоку те же архитектуры становятся инвариантными/эквивариантными без ухудшения равномерной погрешности благодаря TW-устойчивости.
\end{itemize}

%%%%%%%%%%%%%%%%%%%%%%%%%%%%%%%%%%%%%%%


\subsection{Идентификация динамических систем через теорему об аппроксимации операторов}

\textbf{Постановка.} Пусть \( K_1, K_2 \subset \mathbb{R}^n \) компактны, \( V \subset C(K_1) \) компактно. Дискретная динамическая система на функциях задаётся непрерывным оператором

\[
\varphi \colon V \to C(K_2), \quad u_{t+1} = \varphi(u_t).
\]

\begin{theorem}[Теорема D, напоминание]
Для любого \( \varepsilon > 0 \) существуют целые числа \( m, M, N \), точки \( x_1, \ldots, x_m \in K_1 \), <<частоты>> выхода \( w_1, \ldots, w_N \in \mathbb{R}^n \) и параметры \( \{\theta_{ij}, \gamma_i\}_{i \leq M, j \leq m} \), \( \{\alpha_{ik}, \beta_k\}_{i \leq M, k \leq N} \) такие, что
\[
\sup_{u \in V} \sup_{y \in K_2} \bigg| \varphi(u)(y) - \sum_{k=1}^{N} \bigg( \sum_{i=1}^{M} \alpha_{ik} g \big( \sum_{j=1}^{m} \theta_{ij} u(x_j) + \gamma_i \big) \bigg) g(w_k \cdot y + \beta_k) \bigg| \leq \varepsilon.
\]
\end{theorem}

\textit{Интерпретация:} двухслойное сепарабельное (<<внешнее произведение>>) разложение: внутренние \( g \) обрабатывают конечное число входных выборок, внешние \( g \) кодируют зависимость выхода от \( y \).


\begin{example}[Перенос (сдвиг/адвекция) на торе]

Пусть \( K_1 = K_2 = \mathbb{T}^n \), \( \varphi(u)(y) = u(y - a) \) (адвекция на \( a \in \mathbb{R}^n \)). Частотный отклик:  
\[
\widehat{\varphi(u)} (\xi) = e^{-i\xi\cdot a} \, \widehat{u}(\xi).
\]

\textbf{Примененим теорему D:}

\begin{itemize}
\item Выберем ширину полосы \( R \) и решётку \( \Lambda_R = \{ \xi \in \mathbb{Z}^n \colon |\xi| \leq R \} \).
\item Внешние атомы: для каждого \( \xi \in \Lambda_R \) приближем \( \cos(\xi \cdot y) \), \( \sin(\xi \cdot y) \) конечными суммами \( g(w_k \cdot y + \beta_k) \) (ТВ 1D вдоль \( t = \xi \cdot y \)).
\item Внутреннее отображение: по выборкам \( u(x_j) \) линейно восстанавливаем (аппроксимируем) низкочастотные коэффициенты Фурье \( \widehat{u}(\xi) \) (через фиксированную квадратуру/ДПФ на сетке), затем умножаем на \( e^{-i\xi\cdot a} \) для получения целевой значения.
\end{itemize}

\textbf{Полученная сеть.} Произведение <<(внутреннее приближение \( \widehat{u}(\xi) \)) $\times$ (внешний тригонометрический атом на частоте \( \xi \))>> реализует оператор. Обучение не требуется, если \( a \) известно, если \( a \) неизвестно, фазы подбираются методом наименьших квадратов по коэффициентам.

Подробнее рассмотрим случай $n = 1$.
\begin{itemize}
\item Выбираем \( \Lambda_R = \{-R, \ldots, R\} \); на сетке \( Y \) вычисляем ДПФ \( \widehat{u}(\xi) \) и целевые значения \( \widehat{\varphi(u)}(\xi) = e^{-i\xi a}\widehat{u}(\xi) \).
\item Внешний слой: подбираем \( \psi_k(y) = g(w_k y + \beta_k) \) и \( q_{k\xi} \) так, чтобы \( \sum_k q_{k\xi} \psi_k(y) \approx \cos(\xi y), \sin(\xi y) \).
\item Внутренний слой: определяем \( h(u) \) по выборкам \( u(x_j) \); подбираем \( B \) так, чтобы \( B h(u) \approx (\widehat{u}(\xi))_{\xi \in \Lambda_R} \).
\item Собираем: \( C = Q \, \text{diag}( e^{-i\xi a} ) \, B \) (разделяем на блоки косинусов/синусов для вещественной арифметики).
\end{itemize}

\textbf{Результат.} Точность ограничена только уровнем усечения и аппроксимацией ridge-атомов, обратное распространение ошибки не используется.
\end{example}

%%%%%%%%%%%%%%%%%%%

\begin{remarks}
Некоторые практические замечания и направления обобщения.

\textbf{Выбор \( g \).} Любая ограниченная не полиномиальная функция (например, сигмоида, $\th$) является TW, ReLU работает в \( L^p \) и на практике также для тригонометрических аппроксимаций.

\textbf{Вычисления.} Калибровка внешнего слоя и подбор \( B \) сводятся к линейному методу наименьших квадратов; для устойчивости к шуму можно добавить регуляризацию Тихонова. Весь конвейер \textit{амортизирован}: как только \( \{\psi_k\} \) и \( B \) зафиксированы, развёртывание представляет собой быструю двухслойную сеть.

\textbf{За пределами периодических областей.} Чётно продолжаем на \([-1, 1]^n\), затем применяем Фейера/Бохнера–Рисса на торе (как в доказательстве) и сужаем обратно.

\textbf{Симметрии.} Инвариантность/эквивариантность обеспечивается проекцией \( C \) путём усреднения по группе (см. предыдущие темы), оценки ошибок при этом сохраняются.
\end{remarks}

Приведём главные результаты лекции.
\begin{itemize}
\item  \textbf{Активации:} \( g \in C(\mathbb{R}) \cap \mathcal{S}'(\mathbb{R}) \) является TW \( \iff \) \( g \) не является полиномом, ограниченные сигмоидальные функции являются TW.

\item  \textbf{Равномерная УА:} аппроксимация выполняется \textit{равномерно по компактным множествам функций} с параметрами, не зависящими от \( f \) (коэффициенты линейны по \( f \)).

\item \textbf{Функционалы и операторы:} любой непрерывный функционал на \( V \subset C(K) \) и непрерывный оператор \( G \colon V \to C(K_2) \) допускают равномерные нейросетевые аппроксимации, зависящие от конечного числа входных выборок и набора выходных \textit{пространственных атомов}.

\item \textbf{Динамика:} рассматриваем систему как оператор \( K \), симметрии сдвига/потока естественным образом учитываются через семейства сдвигов/растяжений внутри нейросетевого базиса, аппроксимируется вся выходная траектория.
\end{itemize}
Основой стал материал статей Чена и Чена \cite{Chen1995, Chen1993, Chen1993a}.
\end{document}